\documentclass[11pt,a4paper]{article}
\usepackage[utf8]{inputenc}
\usepackage[vietnamese]{babel}
\usepackage{amsmath,amssymb}
\usepackage{geometry}
\geometry{a4paper, total={180mm,260mm}, left=15mm, right=15mm, top=18mm, bottom=18mm}

\usepackage{tikz}
\usetikzlibrary{shapes.geometric, arrows.meta, positioning, calc}
\tikzset{arr/.style={-{Stealth[scale=0.8]}}}
\usepackage{tcolorbox}
\tcbuselibrary{skins,breakable}
\usepackage{fontawesome}
\usepackage{fancyhdr}
\usepackage{enumitem}

\pagestyle{fancy}
\fancyhf{}
\renewcommand{\headrulewidth}{0.5pt}
\renewcommand{\footrulewidth}{0.5pt}
\fancyhead[L]{\small\bfseries\color{violet!85!black}\faBook\ SỔ TAY KIẾN THỨC TOÁN 2}
\fancyhead[R]{\small\bfseries\color{violet!85!black}Lớp Toán Thầy Lĩnh -- loptoanthaylinh.vn}
\fancyfoot[C]{\small\color{gray}-- Sổ tay kiến thức: Bài 06. Luyện tập chung --}

\setlength{\parindent}{0pt}
\setlength{\parskip}{0pt}

\begin{document}

% --- BANNER BÀI HỌC ---
\begin{tcolorbox}[
  enhanced,
  colback=violet!8!white,
  colframe=violet!70!black,
  arc=3mm, boxrule=1.2pt,
  left=3mm, right=3mm, top=2mm, bottom=2mm,
  halign=center
]
  {\LARGE\bfseries\color{violet!90!black}\faBookmark\ BÀI 06: LUYỆN TẬP CHUNG}\\[1mm]
  {\itshape\footnotesize\color{black!75} Đại tổng kết kiến thức Chủ đề 1: Số đến 100, tia số, phép tính và hơn kém nhau trên 1 trang A4}
\end{tcolorbox}

\vspace{-2mm}

\noindent
\begin{minipage}[t]{0.49\textwidth}\vspace{0pt}
\begin{tcolorbox}[
  enhanced,
  colback=violet!4!white,
  colframe=violet!75!black,
  arc=2mm, boxrule=0.8pt,
  left=2mm, right=2mm, top=1.8mm, bottom=1.8mm,
  title={\footnotesize\bfseries\color{white}\faKey\ 1. Trụ cột kiến thức cốt lõi}
]
{\footnotesize
\begin{itemize}\setlength{\itemsep}{1pt}\setlength{\leftskip}{-2.5mm}
  \item \textbf{Số học:} Đọc, viết số đến 100; phân tích $a0 + b$; so sánh số chục trước rồi đến đơn vị.
  \item \textbf{Tia số:} Chiều tăng từ trái qua phải; số liền trước bớt 1, số liền sau thêm 1.
  \item \textbf{Thành phần phép tính:}
  \begin{itemize}\setlength{\itemsep}{0.5pt}\setlength{\leftskip}{-2.5mm}
    \item Phép cộng: $\text{Số hạng} + \text{Số hạng} = \text{Tổng}$.
    \item Phép trừ: $\text{Số bị trừ} - \text{Số trừ} = \text{Hiệu}$.
  \end{itemize}
  \item \textbf{Bài toán chênh lệch:} Hơn/kém nhau bao nhiêu $\implies$ Luôn lấy $\text{Số lớn} - \text{Số bé}$.
\end{itemize}
}
\end{tcolorbox}
\end{minipage}
\hfill
\begin{minipage}[t]{0.49\textwidth}\vspace{0pt}
\begin{tcolorbox}[
  enhanced,
  colback=teal!4!white,
  colframe=teal!75!black,
  arc=2mm, boxrule=0.8pt,
  left=2mm, right=2mm, top=1.8mm, bottom=1.8mm,
  title={\footnotesize\bfseries\color{white}\faSliders\ 2. Kĩ năng tổng hợp}
]
{\footnotesize
\begin{itemize}\setlength{\itemsep}{1pt}\setlength{\leftskip}{-2.5mm}
  \item \textbf{So sánh hai vế biểu thức:}
  \begin{itemize}\setlength{\itemsep}{0.5pt}\setlength{\leftskip}{-2.5mm}
    \item Bước 1: Tính kết quả của từng vế.
    \item Bước 2: So sánh hai kết quả tìm được.
  \end{itemize}
  \item \textbf{Tìm chữ số thích hợp:} Dựa vào hàng đơn vị trước rồi suy ra hàng chục.
  \item \textbf{Giải toán 1 bước:} Đọc kĩ bài toán để xác định phép tính phù hợp (gộp lại $\to$ cộng; bớt đi hoặc tìm chênh lệch $\to$ trừ).
\end{itemize}
}
\end{tcolorbox}
\end{minipage}

\vspace{-1.5mm}
\begin{tcolorbox}[
  enhanced,
  colback=orange!4!white,
  colframe=orange!80!black,
  arc=2mm, boxrule=0.8pt,
  left=2mm, right=2mm, top=1.5mm, bottom=1.5mm,
  title={\footnotesize\bfseries\color{white}\faSitemap\ 3. Sơ đồ tư duy 4 nhánh: Tổng kết Chủ đề 1 (Ôn tập và bổ sung)}
]
\centering
\begin{tikzpicture}[
  node distance=0.8cm,
  box/.style={rectangle, rounded corners=2pt, draw=violet!70!black, fill=violet!10, text width=2.4cm, align=center, font=\scriptsize\bfseries, inner sep=3pt},
  branch/.style={rectangle, rounded corners=2pt, draw=#1!80!black, fill=#1!10, text width=2.3cm, align=center, font=\scriptsize, inner sep=2.5pt},
  arr/.style={-{Stealth[scale=0.8]}, thick, #1!80!black}
]
  \node[box] (root) {CHỦ ĐỀ 1\\ÔN TẬP \& BỔ SUNG};
  \node[branch=teal, above left=0.3cm and 0.4cm of root] (b1) {\textbf{SỐ ĐẾN 100}\\Đọc, viết, so sánh\\Cấu tạo: $a0 + b$};
  \node[branch=orange, above right=0.3cm and 0.4cm of root] (b2) {\textbf{TIA SỐ}\\Gốc 0, tăng dần\\Trước (-1) / Sau (+1)};
  \node[branch=green, below left=0.3cm and 0.4cm of root] (b3) {\textbf{THÀNH PHẦN}\\Số hạng - Tổng\\Số bị trừ - Số trừ - Hiệu};
  \node[branch=purple, below right=0.3cm and 0.4cm of root] (b4) {\textbf{HƠN KÉM NHAU}\\Phần chênh lệch\\$=$ Số lớn $-$ Số bé};

  \draw[arr=teal] (root) -- (b1);
  \draw[arr=orange] (root) -- (b2);
  \draw[arr=green] (root) -- (b3);
  \draw[arr=purple] (root) -- (b4);
\end{tikzpicture}
\end{tcolorbox}

\vspace{-1.5mm}
\noindent
\begin{minipage}[t]{0.49\textwidth}\vspace{0pt}
\begin{tcolorbox}[
  enhanced,
  colback=green!4!white,
  colframe=green!60!black,
  arc=2mm, boxrule=0.8pt,
  left=2mm, right=2mm, top=1.8mm, bottom=1.8mm,
  title={\footnotesize\bfseries\color{white}\faCheckCircle\ 4. Ví dụ thực hành chuẩn mực}
]
{\footnotesize
\textbf{1. Điền dấu ($>, <, =$):}\\
So sánh: $35 + 14 \dots 50$.\\
Vế trái: $35 + 14 = 49$.\\
Vì $49 < 50$ nên: $35 + 14 < 50$.
\textbf{2. Bài toán lời văn:}\\
Lớp 2A có $35$ học sinh, lớp 2B có $32$ học sinh. Hỏi lớp 2A nhiều hơn lớp 2B bao nhiêu học sinh?\\
\textit{Bài giải:}\\
Lớp 2A nhiều hơn lớp 2B số học sinh là:
$$35 - 32 = 3\text{ (học sinh)}$$
}
\end{tcolorbox}
\end{minipage}
\hfill
\begin{minipage}[t]{0.49\textwidth}\vspace{0pt}
\begin{tcolorbox}[
  enhanced,
  colback=red!4!white,
  colframe=red!75!black,
  arc=2mm, boxrule=0.8pt,
  left=2mm, right=2mm, top=1.8mm, bottom=1.8mm,
  title={\footnotesize\bfseries\color{white}\faExclamationTriangle\ 5. Mẹo tránh bẫy trắc nghiệm}
]
{\footnotesize
\begin{itemize}\setlength{\itemsep}{1pt}\setlength{\leftskip}{-2.5mm}
  \item \textbf{Bẫy so sánh biểu thức chưa tính:}\\
  Tuyệt đối không đoán mò, phải tính giá trị hai vế ra nháp rồi mới điền dấu.
  \item \textbf{Bẫy số liền trước số tròn chục:}\\
  Số liền trước của $50$ là $49$, số liền trước của $70$ là $69$.
  \item Luôn viết đáp số kèm theo đơn vị đo trong ngoặc.
\end{itemize}
}
\end{tcolorbox}
\end{minipage}


\end{document}
